\documentclass[10pt,a4paper]{amsart}
\usepackage{amsmath,amssymb,amsthm,graphicx}
% Retained minimal preamble for the Stacks Project chapter sites-cohomology.tex (Cohomology on Sites).
% Verbatim excerpt of the chapter preamble: Xy-pic setup (native preamble lines 14-18) and the native
% theorem declarations with the mathematical macros (native preamble lines 154-217). Copied mechanically;
% nothing added, removed or redefined. The wrapper class and the amsmath/amssymb/amsthm/graphicx base
% packages are supplied by the assembler, so the class/externaldocument part of the native preamble is omitted.
% For commutative diagrams we use Xy-pic
\usepackage[all]{xy}

% We use 2cell for 2-commutative diagrams.
\xyoption{2cell}
% Theorem environments.
%
\theoremstyle{plain}
\newtheorem{theorem}[subsection]{Theorem}
\newtheorem{proposition}[subsection]{Proposition}
\newtheorem{lemma}[subsection]{Lemma}

\theoremstyle{definition}
\newtheorem{definition}[subsection]{Definition}
\newtheorem{example}[subsection]{Example}
\newtheorem{exercise}[subsection]{Exercise}
\newtheorem{situation}[subsection]{Situation}

\theoremstyle{remark}
\newtheorem{remark}[subsection]{Remark}
\newtheorem{remarks}[subsection]{Remarks}

\numberwithin{equation}{subsection}

% Macros
%
\def\lim{\mathop{\mathrm{lim}}\nolimits}
\def\colim{\mathop{\mathrm{colim}}\nolimits}
\def\Spec{\mathop{\mathrm{Spec}}}
\def\Hom{\mathop{\mathrm{Hom}}\nolimits}
\def\Ext{\mathop{\mathrm{Ext}}\nolimits}
\def\SheafHom{\mathop{\mathcal{H}\!\mathit{om}}\nolimits}
\def\SheafExt{\mathop{\mathcal{E}\!\mathit{xt}}\nolimits}
\def\Sch{\mathit{Sch}}
\def\Mor{\mathop{\mathrm{Mor}}\nolimits}
\def\Ob{\mathop{\mathrm{Ob}}\nolimits}
\def\Sh{\mathop{\mathit{Sh}}\nolimits}
\def\NL{\mathop{N\!L}\nolimits}
\def\CH{\mathop{\mathrm{CH}}\nolimits}
\def\proetale{{pro\text{-}\acute{e}tale}}
\def\etale{{\acute{e}tale}}
\def\QCoh{\mathit{QCoh}}
\def\Ker{\mathop{\mathrm{Ker}}}
\def\Im{\mathop{\mathrm{Im}}}
\def\Coker{\mathop{\mathrm{Coker}}}
\def\Coim{\mathop{\mathrm{Coim}}}

% Boxtimes
%
\DeclareMathSymbol{\boxtimes}{\mathbin}{AMSa}{"02}

%
% Macros for moduli stacks/spaces
%
\def\QCohstack{\mathcal{QC}\!\mathit{oh}}
\def\Cohstack{\mathcal{C}\!\mathit{oh}}
\def\Spacesstack{\mathcal{S}\!\mathit{paces}}
\def\Quotfunctor{\mathrm{Quot}}
\def\Hilbfunctor{\mathrm{Hilb}}
\def\Curvesstack{\mathcal{C}\!\mathit{urves}}
\def\Polarizedstack{\mathcal{P}\!\mathit{olarized}}
\def\Complexesstack{\mathcal{C}\!\mathit{omplexes}}
% \Pic is the operator that assigns to X its picard group, usage \Pic(X)
% \Picardstack_{X/B} denotes the Picard stack of X over B
% \Picardfunctor_{X/B} denotes the Picard functor of X over B
\def\Pic{\mathop{\mathrm{Pic}}\nolimits}
\def\Picardstack{\mathcal{P}\!\mathit{ic}}
\def\Picardfunctor{\mathrm{Pic}}
\def\Deformationcategory{\mathcal{D}\!\mathit{ef}}
\usepackage{cleveref}
\crefname{originallemma}{original lemma}{original lemmas}
\Crefname{originallemma}{Original lemma}{Original lemmas}
\makeatletter
\newlabel{lemma-adjoints-push-pull-compatibility}{{0FP1}{source}{}{sourcepointer-1}{}}
\newlabel{lemma-adjoints-push-pull-compatibility@cref}{{[originallemma][0][]0FP1}{[1][0][]source}{}{}{}}
\makeatother
\crefname{originallemma}{original lemma}{original lemmas}
\Crefname{originallemma}{Original lemma}{Original lemmas}
\makeatletter
\newlabel{lemma-tensor-pull-compatibility}{{0FP0}{source}{}{sourcepointer-2}{}}
\newlabel{lemma-tensor-pull-compatibility@cref}{{[originallemma][0][]0FP0}{[1][0][]source}{}{}{}}
\makeatother
\crefname{originalremark}{original remark}{original remarks}
\Crefname{originalremark}{Original remark}{Original remarks}
\makeatletter
\newlabel{remark-cup-product}{{0B68}{source}{}{sourcepointer-3}{}}
\newlabel{remark-cup-product@cref}{{[originalremark][0][]0B68}{[1][0][]source}{}{}{}}
\makeatother
\title{The cup product on derived pushforwards agrees with the naive cup product on total tensor complexes}
\author{The Stacks Project Authors}
\date{stacks-project:sites-cohomology.tex:local-snapshot-20261011; GNU FDL 1.2 or later}
\begin{document}
\maketitle
\noindent\textit{Editorial note.} All delivered mathematical text is the author's own, copied from the pinned Stacks Project chapter file sites-cohomology.tex; no statement or proof step was rewritten, reordered or removed.\par\medskip
\noindent\textit{Editorial note.} The counted proof is the author's complete argument: the proof is written out in full: both composites around the displayed diagram are identified as adjoints of explicit maps, normalised by the adjoint push-pull compatibility lemma, and the final diagram chase names the reason for the commutativity of every polygon, ending with the observation that the outer composite is the pair of maps in question.\par\medskip
\noindent\textit{Editorial note.} The label lemma-cup-compatible-with-naive names the counted statement, and the proof environment immediately following it in the source is delivered verbatim as its proof.\par\medskip
\noindent\textit{Editorial note.} The statement is a displayed commutative diagram in the author's own notation; it is delivered verbatim, including the references naming the cup product and the naive cup product.\par\medskip
\noindent\textit{Editorial note.} The preamble is the chapter's own retained minimal preamble (sources/stacks-sites-cohomology/retained-preamble.tex), a verbatim excerpt of the native preamble: the Xy-pic setup and the native theorem declarations with the mathematical macros.\par\medskip
\noindent\textit{Allowed standard prerequisites (editorial).}\begin{itemize}
\item The naive cup product on total tensor complexes of modules on a ringed site. Same chapter, the naive cup product Tot(K\textasciicircum{}bullet tensor M\textasciicircum{}bullet) -> Tot(K\textasciicircum{}bullet tensor\textasciicircum{}L M\textasciicircum{}bullet) and its adjoint description. Supplies the vertical maps of the displayed diagram and the two naive cup products compared in it.
\item Cup product on derived pushforwards via the derived tensor product. Same chapter, the cup product defined by f\_*K tensor\textasciicircum{}L f\_*M -> Rf\_*(K tensor\textasciicircum{}L M) as the adjoint of the canonical map. Supplies the upper and right-hand maps of the diagram whose commutativity is asserted.
\item Adjunction between Lf\textasciicircum{}* and Rf\_* and the comparison Lf\textasciicircum{}* -> f\textasciicircum{}*. Same chapter, the adjunction isomorphisms for the morphism of ringed topoi f and the functorial comparison between the derived and underived pullbacks. Supplies the identification of both composites as the same adjoint map, which is how the two directions around the diagram are shown to agree.
\end{itemize}\medskip
\ifdefined\hypertarget\hypertarget{sourcepointer-1}{}\fi
\noindent\textit{Source reference.} Original lemma 0FP1: sources/stacks-sites-cohomology/pinned-original-source.tex, source lines 4129--4129. Prerequisite pointer invoked at proof line(s) 8306; it names a standard fact of the same chapter that the counted argument uses, and every step of the counted proof built on it is retained verbatim in the delivered body.\par\medskip
\ifdefined\hypertarget\hypertarget{sourcepointer-2}{}\fi
\noindent\textit{Source reference.} Original lemma 0FP0: sources/stacks-sites-cohomology/pinned-original-source.tex, source lines 3849--3849. Prerequisite pointer invoked at proof line(s) 8407; it names a standard fact of the same chapter that the counted argument uses, and every step of the counted proof built on it is retained verbatim in the delivered body.\par\medskip
\ifdefined\hypertarget\hypertarget{sourcepointer-3}{}\fi
\noindent\textit{Source reference.} Original remark 0B68: sources/stacks-sites-cohomology/pinned-original-source.tex, source lines 4193--4193. Prerequisite pointer invoked at proof line(s) 8254; it names a standard fact of the same chapter that the counted argument uses, and every step of the counted proof built on it is retained verbatim in the delivered body.\par\medskip
\begin{lemma}
\label{lemma-cup-compatible-with-naive}
In the situation above the following diagram commutes
$$
\xymatrix{
f_*\mathcal{K}^\bullet
\otimes_{\mathcal{O}_\mathcal{D}}^\mathbf{L}
f_*\mathcal{M}^\bullet \ar[r] \ar[d]
&
Rf_*\mathcal{K}^\bullet
\otimes_{\mathcal{O}_\mathcal{D}}^\mathbf{L}
Rf_*\mathcal{M}^\bullet \ar[d]^{\text{Remark \ref{remark-cup-product}}} \\
\text{Tot}(
f_*\mathcal{K}^\bullet
\otimes_{\mathcal{O}_\mathcal{D}}
f_*\mathcal{M}^\bullet) \ar[d]_{\text{naive cup product}} &
Rf_*(\mathcal{K}^\bullet
\otimes_{\mathcal{O}_\mathcal{C}}^\mathbf{L}
\mathcal{M}^\bullet) \ar[d] \\
f_*\text{Tot}(\mathcal{K}^\bullet
\otimes_{\mathcal{O}_\mathcal{C}}
\mathcal{M}^\bullet) \ar[r] &
Rf_*\text{Tot}(\mathcal{K}^\bullet
\otimes_{\mathcal{O}_\mathcal{C}}
\mathcal{M}^\bullet)
}
$$
\end{lemma}

\begin{proof}
By the construction in Remark \ref{remark-cup-product} we see that
going around the diagram clockwise the map
$$
f_*\mathcal{K}^\bullet
\otimes_{\mathcal{O}_\mathcal{D}}^\mathbf{L}
f_*\mathcal{M}^\bullet 
\longrightarrow
Rf_*\text{Tot}(\mathcal{K}^\bullet
\otimes_{\mathcal{O}_\mathcal{C}}
\mathcal{M}^\bullet)
$$
is adjoint to the map
\begin{align*}
Lf^*(f_*\mathcal{K}^\bullet
\otimes_{\mathcal{O}_\mathcal{D}}^\mathbf{L}
f_*\mathcal{M}^\bullet)
& =
Lf^*f_*\mathcal{K}^\bullet
\otimes_{\mathcal{O}_\mathcal{D}}^\mathbf{L}
Lf^*f_*\mathcal{M}^\bullet \\
& \to
Lf^*Rf_*\mathcal{K}^\bullet
\otimes_{\mathcal{O}_\mathcal{D}}^\mathbf{L}
Lf^*Rf_*\mathcal{M}^\bullet \\
& \to
\mathcal{K}^\bullet
\otimes_{\mathcal{O}_\mathcal{D}}^\mathbf{L}
\mathcal{M}^\bullet \\
& \to
\text{Tot}(\mathcal{K}^\bullet
\otimes_{\mathcal{O}_\mathcal{C}}
\mathcal{M}^\bullet)
\end{align*}
By Lemma \ref{lemma-adjoints-push-pull-compatibility} this is also equal to
\begin{align*}
Lf^*(f_*\mathcal{K}^\bullet
\otimes_{\mathcal{O}_\mathcal{D}}^\mathbf{L}
f_*\mathcal{M}^\bullet)
& =
Lf^*f_*\mathcal{K}^\bullet
\otimes_{\mathcal{O}_\mathcal{D}}^\mathbf{L}
Lf^*f_*\mathcal{M}^\bullet \\
& \to
f^*f_*\mathcal{K}^\bullet
\otimes_{\mathcal{O}_\mathcal{D}}^\mathbf{L}
f^*f_*\mathcal{M}^\bullet \\
& \to
\mathcal{K}^\bullet
\otimes_{\mathcal{O}_\mathcal{D}}^\mathbf{L}
\mathcal{M}^\bullet \\
& \to
\text{Tot}(\mathcal{K}^\bullet
\otimes_{\mathcal{O}_\mathcal{C}}
\mathcal{M}^\bullet)
\end{align*}
Going around anti-clockwise we obtain the map adjoint to the map
\begin{align*}
Lf^*(f_*\mathcal{K}^\bullet
\otimes_{\mathcal{O}_\mathcal{D}}^\mathbf{L}
f_*\mathcal{M}^\bullet)
& \to
Lf^*\text{Tot}(
f_*\mathcal{K}^\bullet
\otimes_{\mathcal{O}_\mathcal{D}}
f_*\mathcal{M}^\bullet) \\
& \to
Lf^*f_*\text{Tot}(\mathcal{K}^\bullet
\otimes_{\mathcal{O}_\mathcal{C}}
\mathcal{M}^\bullet) \\
& \to
Lf^*Rf_*\text{Tot}(\mathcal{K}^\bullet
\otimes_{\mathcal{O}_\mathcal{C}}
\mathcal{M}^\bullet) \\
& \to
\text{Tot}(\mathcal{K}^\bullet
\otimes_{\mathcal{O}_\mathcal{C}}
\mathcal{M}^\bullet)
\end{align*}
By Lemma \ref{lemma-adjoints-push-pull-compatibility} this is also equal to
\begin{align*}
Lf^*(f_*\mathcal{K}^\bullet
\otimes_{\mathcal{O}_\mathcal{D}}^\mathbf{L}
f_*\mathcal{M}^\bullet)
& \to
Lf^*\text{Tot}(
f_*\mathcal{K}^\bullet
\otimes_{\mathcal{O}_\mathcal{D}}
f_*\mathcal{M}^\bullet) \\
& \to
Lf^*f_*\text{Tot}(\mathcal{K}^\bullet
\otimes_{\mathcal{O}_\mathcal{C}}
\mathcal{M}^\bullet) \\
& \to
f^*f_*\text{Tot}(\mathcal{K}^\bullet
\otimes_{\mathcal{O}_\mathcal{C}}
\mathcal{M}^\bullet) \\
& \to
\text{Tot}(\mathcal{K}^\bullet
\otimes_{\mathcal{O}_\mathcal{C}}
\mathcal{M}^\bullet)
\end{align*}
Now the proof is finished by a contemplation of the diagram
$$
\xymatrix{
Lf^*(f_*\mathcal{K}^\bullet
\otimes_{\mathcal{O}_\mathcal{D}}^\mathbf{L}
f_*\mathcal{M}^\bullet) \ar[d] \ar[rr] & &
Lf^*f_*\mathcal{K}^\bullet \otimes_{\mathcal{O}_\mathcal{C}}^\mathbf{L}
Lf^*f_*\mathcal{M}^\bullet \ar[d] \\
Lf^*\text{Tot}(
f_*\mathcal{K}^\bullet
\otimes_{\mathcal{O}_\mathcal{D}}
f_*\mathcal{M}^\bullet) \ar[d]_{naive} \ar[r] &
f^*\text{Tot}(
f_*\mathcal{K}^\bullet
\otimes_{\mathcal{O}_\mathcal{D}}
f_*\mathcal{M}^\bullet) \ar[ldd]^{naive} \ar[dd] &
f^*f_*\mathcal{K}^\bullet \otimes_{\mathcal{O}_\mathcal{C}}^\mathbf{L}
f^*f_*\mathcal{M}^\bullet \ar[dd] \ar[ldd] \\
Lf^*f_*\text{Tot}(\mathcal{K}^\bullet
\otimes_{\mathcal{O}_\mathcal{C}}
\mathcal{M}^\bullet) \ar[d] \\
f^*f_*\text{Tot}(\mathcal{K}^\bullet \otimes_{\mathcal{O}_\mathcal{C}}
\mathcal{M}^\bullet) \ar[rd] &
\text{Tot}(f^*f_*\mathcal{K}^\bullet \otimes_{\mathcal{O}_\mathcal{C}}
f^*f_*\mathcal{M}^\bullet) \ar[d] &
\mathcal{K}^\bullet \otimes_{\mathcal{O}_\mathcal{C}}^\mathbf{L}
\mathcal{M}^\bullet \ar[ld] \\
& \text{Tot}(\mathcal{K}^\bullet
\otimes_{\mathcal{O}_\mathcal{C}}
\mathcal{M}^\bullet)
}
$$
All of the polygons in this diagram commute. The top one commutes
by Lemma \ref{lemma-tensor-pull-compatibility}.
The square with the two naive cup products commutes because
$Lf^* \to f^*$ is functorial in the complex of modules.
Similarly with the square involving the two maps
$\mathcal{A}^\bullet \otimes^\mathbf{L} \mathcal{B}^\bullet \to
\text{Tot}(\mathcal{A}^\bullet \otimes \mathcal{B}^\bullet)$.
Finally, the commutativity of the remaining square
is true on the level of complexes and may be viewed as the
definition of the naive cup product (by the adjointness
of $f^*$ and $f_*$). The proof is finished because
going around the diagram on the outside are the two maps
given above.
\end{proof}

\end{document}
